% ============================================================
% Zentrales Layout fuer das Chemie-Skript
% Absichtlich einfach gehalten: keine absoluten Positionen.
% Compiler in Overleaf: pdfLaTeX
% ============================================================

% --- Pakete --------------------------------------------------
\usepackage[T1]{fontenc}
\usepackage[utf8]{inputenc}
\usepackage{xcolor}       % Farben fuer Ueberschriften und Boxen
\usepackage[a4paper,margin=20mm,top=25mm,bottom=22mm,
            headheight=20pt,headsep=8mm,footskip=12mm]{geometry} % Seitenraender
\usepackage{graphicx}     % Abbildungen einbinden
\usepackage{amsmath}      % Mathematische Formeln (z.B. Brueche)
\usepackage{needspace}    % Aufgaben samt nachfolgender Abbildung zusammenhalten
\usepackage{fancyhdr}     % Kopf- und Fusszeile
\usepackage{titlesec}     % Normale LaTeX-Ueberschriften gestalten
\usepackage{microtype}    % Typografisch bessere Buchstaben- und Wortabstaende
\usepackage{parskip}      % Absatzabstand statt Absatzeinzug
\usepackage{tcolorbox}    % Blaue Infoboxen
\usepackage{tabularray}   % Tabellen mit gut kontrollierbaren Zellen
% Keine automatisch erzeugten 'Table 1'-Beschriftungen an langen Arbeitsblatt-Tabellen.
\SetTblrTemplate{caption}{empty}
\SetTblrTemplate{capcont}{empty}
\SetTblrTemplate{contfoot}{empty}
\usepackage{wrapfig}      % Textumfluss bei Bildern (in einigen Kapiteln)
\usepackage{ragged2e}     % Verbesserter Flattersatz mit Silbentrennung
\usepackage{enumitem}     % Anpassbare Aufzählungen und Nummerierungen
\usepackage[outline]{contour} % Umrandung einzelner Buchstaben
\usepackage[british]{babel} % Britisches Englisch und Silbentrennung
\usepackage[hidelinks]{hyperref} % Klickbare Links im PDF
\usepackage{pgfplots} % Diagramme direkt in LaTeX
\pgfplotsset{compat=1.18}
\usetikzlibrary{patterns}
\usepackage{caption}
\usepackage{chemfig}
\usepackage[version=4]{mhchem}

% --- Schrift und Farben --------------------------------------
\usepackage{lmodern}

\definecolor{DeepBlue}{HTML}{233B53}
\definecolor{InfoBlue}{HTML}{ECF9FF}
\definecolor{TableGrid}{HTML}{CCD3D9}

% Absichtlich normale LaTeX-Trennung und Textsatz verwenden.
% Keine manuellen globalen Zeilenhoehen oder Umbruchverbote.

% --- Zeilenabstand -------------------------------------------
\linespread{1.08} % Etwas lockererer Zeilenabstand
\setlength{\parindent}{0pt} % Keine Einrückung zu Beginn eines Absatzes
\setlength{\parskip}{0.4em} % Kleiner Abstand zwischen Absätzen

% --- Kopf- und Fusszeile -------------------------------------
\pagestyle{fancy}
\fancyhf{}
\fancyhead[L]{\includegraphics[height=12pt]{images/gkg-logo.png}\hspace{0.4em}\footnotesize\bfseries Chemistry GKG (Nl)}
\fancyhead[C]{\footnotesize\bfseries\CurrentYear}
\fancyhead[R]{\footnotesize\bfseries\ScriptTitle}
\fancyfoot[L]{\scriptsize\copyright\ GKG}
\fancyfoot[C]{\footnotesize\thepage}
\renewcommand{\headrulewidth}{0.4pt}
\renewcommand{\footrulewidth}{0pt}

% --- Titel ---------------------------------------------------
% Weisse Schrift mit dunkelblauer Buchstabenkontur
\newcommand{\OutlinedChapterTitle}[1]{%
  \begin{center}
    \contourlength{0.8pt}
    {\sffamily\itshape\fontsize{28}{28}\selectfont
      \textcolor{white}{\contour{DeepBlue}{#1}}}
      \vspace{10pt}
  \end{center}
}

% Bemerkung am unteren rechten Seitenrand
\newcommand{\ObjectivesNote}[1]{%
  \par\vfill
  \hfill{\color{DeepBlue}\rule{20mm}{0.8pt}}\par
  \smallskip
  \hfill{\footnotesize\bfseries #1}\par
}

% --- Ueberschriften ------------------------------------------
% Regulare LaTeX-Gliederung; spaeter auch \section{...} etc. moeglich.
\titleformat{\section}{\Large\bfseries\color{DeepBlue}}{\thesection}{0.6em}{}
\titleformat{\subsection}{\large\bfseries\color{DeepBlue}}{\thesubsection}{0.6em}{}
\titleformat{\subsubsection}{\normalsize\bfseries\color{DeepBlue}}{\thesubsubsection}{0.6em}{}
\titleformat{\paragraph}[block]{\normalfont\normalsize\bfseries}
{}{0pt}{}
\titlespacing*{\section}{0pt}{2.2ex plus .6ex}{1ex}
\titlespacing*{\subsection}{0pt}{1.6ex plus .4ex}{.7ex}
\titlespacing*{\subsubsection}{0pt}{1.2ex plus .3ex}{.5ex}
\titlespacing*{\paragraph}{0pt}{0.8ex}{0ex}

% Kompatibilitaet mit den bereits erstellten Kapiteln.
% Die Nummern stehen dort noch im Text, deshalb Sternvarianten.
\newcommand{\HOne}[1]{\section{#1}}
\newcommand{\HTwo}[1]{\subsection{#1}}
\newcommand{\HThree}[1]{\subsubsection*{#1}}
\newcommand{\HFour}[1]{\paragraph*{#1}}

% --- Infoboxen -------------------------------------------------
% Einheitlicher Abstand vor und nach Infoboxen
\newlength{\BoxSpacing}
\setlength{\BoxSpacing}{2ex}
% Infobox
% Gleiche Innenabstaende oben und unten, keine negativen Abstaende.
\newtcolorbox{InfoBox}{%
  colback=InfoBlue,colframe=DeepBlue,
  boxrule=0.6pt,arc=0pt,
  left=6pt,right=6pt,top=6pt,bottom=6pt,
  boxsep=0pt,before skip=\BoxSpacing,after skip=\BoxSpacing,
  before upper={\RaggedRight\setlength{\parskip}{0pt}}
  }
\newcommand{\BoxHeading}[1]{%
  {\bfseries\color{DeepBlue}#1}\par\smallskip}

% Zweispaltige Infobox
% Direkter Tabelleninhalt statt verschachtelter \parbox-Bloecke:
% rowsep bestimmt den symmetrischen Abstand zur oberen/unteren Kante.

\newcommand{\CompareBox}[4]{%
  \begin{tcolorbox}[
    colback=InfoBlue,colframe=DeepBlue,
    boxrule=0.6pt,arc=0pt,
    left=0pt,right=0pt,top=0pt,bottom=0pt,
    boxsep=0pt,before skip=\BoxSpacing,after skip=\BoxSpacing
  ]
    \begin{tblr}{
      width=\linewidth,
      colspec={X[l,t] X[l,t]},
      colsep=6pt,rowsep=6pt,
      vline{2}={0.6pt,DeepBlue}
    }
      {\bfseries\color{DeepBlue}#1}\par\smallskip #2 &
      {\bfseries\color{DeepBlue}#3}\par\smallskip #4
    \end{tblr}
  \end{tcolorbox}
}


% --- Standardtabellen ----------------------------------------
% Einheitliches Design, Spaltenaufteilung bleibt individuell.

\NewTblrEnviron{StandardTable}

\SetTblrInner[StandardTable]{
  width=\linewidth,
  colsep=4pt,
  rowsep=3pt,
  % Normale Tabellenlinien
  hlines={0.8pt,TableGrid},
  vlines={0.8pt,TableGrid},
  % Kopfzeile
  row{1}={bg=InfoBlue,font=\bfseries\color{DeepBlue}},
  hline{1,2}={1pt,DeepBlue},
  vline{1-Z}={1}{1pt,DeepBlue}
}

% Abstände der Standardtabellen
\newlength{\TableSpacing}
\setlength{\TableSpacing}{1ex}

% Abstand vor und nach jeder StandardTable
\AddToHook{env/StandardTable/before}{\par\addvspace{\TableSpacing}}
\AddToHook{env/StandardTable/after}{\par\addvspace{\TableSpacing}}

% --- Aufgaben ------------------------------------------------

% Kapitelnummer und automatischer Aufgabenzähler
\newcounter{exercise}
\newcommand{\ExerciseChapter}{\ScriptNumber}
\renewcommand{\theexercise}{\ExerciseChapter-\arabic{exercise}}

\newcommand{\Exercise}[2][]{%
  \refstepcounter{exercise}
  % Optionalen Bezeichner fuer Referenzen speichern
  \if\relax\detokenize{#1}\relax
  \else
  \label{ex:#1}%
  \fi
  \par\smallskip
  \begin{tblr}{
    width=\linewidth,
    colspec={Q[l,m,wd=16mm] X[l,m]},
    colsep=5pt,rowsep=3pt,
    hlines={0.6pt,DeepBlue},vlines={0.6pt,DeepBlue},
    column{1}={bg=InfoBlue,font=\bfseries\color{DeepBlue}},
    column{2}={font=\bfseries}
  }
    Ex \theexercise & #2
  \end{tblr}\par\smallskip
}


% Teilaufgaben
% Automatische Nummerierung a), b), c), ...
\newlist{ExerciseParts}{enumerate}{1}

\setlist[ExerciseParts]{
  label=\alph*),
  leftmargin=*,
  topsep=0.5ex,
  itemsep=0.5ex,
  parsep=0pt
}


% Freier Platz fuer handschriftliche Antworten.
\newcommand{\ResponseGap}[1][38pt]{\par\vspace{#1}}


% Loesungen erhalten die aktuelle Exercise-Nummer aus dem jeweiligen Label.
% Ein Listeneintrag sorgt dafuer, dass auch alle Folgezeilen und Absaetze
% linksbuendig an derselben Textkante beginnen (haengende Nummer).
% Syntax: \Solution{exercise-label}{Text der Loesung ...}
\newcommand{\Solution}[2]{%
  \begin{list}{}{%
    \setlength{\leftmargin}{20mm}%
    \setlength{\labelwidth}{16mm}%
    \setlength{\labelsep}{4mm}%
    \setlength{\rightmargin}{0pt}%
    \setlength{\topsep}{0.6\baselineskip}%
    \setlength{\itemsep}{0pt}%
    \setlength{\parsep}{0.3em}%
    \setlength{\partopsep}{0pt}%
  }
    \item[\textbf{\ref{ex:#1}}] #2
  \end{list}%
}

% Einfache Ziele-Liste: gleicher linksbuendiger Satz fuer alle Bullet Points.
\newcommand{\Objective}[1]{\par\noindent\hangindent=1.6em\hangafter=1 \textbullet\enspace #1\par}

% --- Zeichenfelder -------------------------------------------
% Weisse Box mit schwarzem Rahmen.
% Breite und Höhe werden individuell festgelegt.

\newtcolorbox{DrawingBox}[2]{%
  width=#1,
  height=#2,
  colback=white,
  colframe=black,
  boxrule=0.8pt,
  arc=0pt,
  boxsep=0pt,
  before skip=0pt,
  after skip=0pt
}
